% jhanswer-zh.tex -- 《线性代数（习题解答）》中文版主文件
% 前置条件：先编译过 book-zh.tex（其会写出中文答案文件 bookans.tex）。
% 编译：在 src/zh/ 目录下（或仓库根目录 `make zh-answers`）：
%   TEXINPUTS="<仓库>/src//:" xelatex -interaction=nonstopmode jhanswer-zh.tex
% 结构照抄英文版 jhanswer.tex；差异：
%   * 载入 zhsetup（字体、名称汉化、第X章/第I节 页眉）与 zhcover（中文封面宏）
%   * 前言为中文
%   * 答案条目超链接指向 book-zh.pdf（英文版为 book.pdf）
%   * 读入 bookans-zh.tex（由 Makefile 从 bookans.tex 生成，节标题前缀已汉化）
\documentclass[twoside]{book}
\usepackage{color,graphicx}
\usepackage{answerjh,linalgjh}
\usepackage{zhsetup}
\usepackage{zhcover}

\includeonly{coverans,symlist}

% 答案条目链接指向中文版主书（answerjh 默认硬编码 book.pdf）
\makeatletter
\renewenvironment{ans}[1]{%
  \begin{answerlist}%
    \ifinsidetopic
      \item[\hyperref{book-zh.pdf}{exercise}{\thesubsection.#1}{#1}\hypertarget{ans.\thesubsection.#1}{}]
    \else
      \item[\hyperref{book-zh.pdf}{exercise}{#1}{#1}\hypertarget{ans.#1}{}]
    \fi }{%
  \end{answerlist}
}
\makeatother

% 读入由 Makefile 生成的汉化版答案汇编（缺省机制同英文版）
\makeatletter
\def\ansfile{bookans-zh}
\makeatother

\begin{document}
% 中文版答案书封面（zhcover 由 jhanswer-zh.tex 在导言区加载）
\thispagestyle{empty}
\setlength{\unitlength}{1in}
\noindent\begin{picture}(0,0)
  \put(-.75,0){\covergraphic[习题解答]}
\end{picture}
\thispagestyle{empty}
% \vfill
%\medskip
\begin{center}
  \textbf{记号}  \\[2ex]
  \begin{tabular}{r|l}
    \( \Re \), \( \Re^+ \), \( \Re^n \) &实数、正实数、实数 $n$ 元组 \\
    \( \N              \),
    \( \C              \)  &自然数 \( \set{0,1,2,\ldots} \)、复数                           \\
    \( (a\,..\,b) \), \( [a\,..\,b] \) &开区间、闭区间   \\
    \( \sequence{\ldots} \)&序列（次序有影响的列表）    \\
    \( h_{i,j} \)          &矩阵 $H$ 的第~\( i \)~行第~\( j \)~列元素 \\
    \( V,W,U \)            &向量空间               \\
    \( \vec{v} \),
    $\zero$, $\zero_V$     &向量、零向量、空间~$V$ 的零向量   \\
    \( \polyspace_n \), \( \matspace_{\nbym{n}{m}} \)
                          &\( n \)~次多项式空间、\( \nbym{n}{m} \) 矩阵的全体      \\
    \( \spanof{S} \)       &集合的张成                   \\
    \( \sequence{B,D} \), \( \vec{\beta},\vec{\delta} \)
                          &基、基向量  \\
    \( \stdbasis_n=\sequence{\vec{e}_1,\,\ldots,\,\vec{e}_n} \)
                          &\( \Re^n \) 的标准基  \\
    \( V\isomorphicto W \) &同构的空间                         \\
    \( M\directsum N \)    &子空间的直和                   \\
    \( h,g \)              &同态（线性映射）                \\
    % \( H,G \)              &matrices                                  \\
    \( t,s \)              &变换（空间到自身的线性映射） \\
    % \( T,S \)              &square matrices                           \\
    \( \rep{\vec{v}}{B} \), \( \rep{h}{B,D} \)
                          &向量、映射的表示    \\
    \( Z_{\nbym{n}{m}} \) 或 \( Z \), \(I_{\nbyn{n}}\) 或 \( I \)  &零矩阵、单位矩阵    \\
    \( \deter{T} \)        &矩阵的行列式       \\
    \( \rangespace{h},\nullspace{h} \)
                           &映射的值域空间、零空间  \\
    \( \genrangespace{h},\gennullspace{h} \)
                           &广义值域空间与广义零空间
  \end{tabular}
\end{center}
\vspace*{\fill}
\begin{center}
  \textbf{希腊字母及其读音}
    \\[1.5ex]
  \newcommand{\pronounced}[1]{\hspace*{.2em}\small\textit{#1}}
  \begin{tabular}{cl@{\hspace*{3em}}cl}
    字符 &\multicolumn{1}{c}{\makebox[-3.5em][r]{名称}}
    &字符  &\multicolumn{1}{c}{\makebox[-3.5em][r]{名称}}  \\
    \hline
     \makebox[1em][l]{\( \alpha  \)} &alpha（阿尔法） \pronounced{AL-fuh}
       &\makebox[1em][l]{\( \nu     \)}  &nu（纽）  \pronounced{NEW}       \\
     \makebox[1em][l]{\( \beta   \)} &beta（贝塔）  \pronounced{BAY-tuh}
       &\makebox[1em][l]{\( \xi  \), \( \Xi \)}  &xi（克西）   \pronounced{KSIGH}    \\
     \makebox[1em][l]{\( \gamma  \), \( \Gamma \)} &gamma（伽马）  \pronounced{GAM-muh}
       &\makebox[1em][l]{\( o \)} &omicron（奥米克戎）  \pronounced{OM-uh-CRON}  \\
     \makebox[1em][l]{\( \delta  \), \( \Delta \)} &delta（德尔塔） \pronounced{DEL-tuh}
       &\makebox[1em][l]{\( \pi \), \( \Pi \)} &pi（派）  \pronounced{PIE}     \\
     \makebox[1em][l]{\( \epsilon\)} &epsilon（艾普西隆）  \pronounced{EP-suh-lon}
       &\makebox[1em][l]{\( \rho \)} &rho（柔）  \pronounced{ROW}    \\
     \makebox[1em][l]{\( \zeta   \)} &zeta（泽塔）   \pronounced{ZAY-tuh}
       &\makebox[1em][l]{\( \sigma  \), \( \Sigma \)} &sigma（西格玛）  \pronounced{SIG-muh}  \\
     \makebox[1em][l]{\( \eta  \)} &eta（艾塔）  \pronounced{AY-tuh}
       &\makebox[1em][l]{\( \tau \)} &tau（陶）  \pronounced{TOW (as in cow)}    \\
     \makebox[1em][l]{\( \theta \), \( \Theta \)} &theta（西塔）  \pronounced{THAY-tuh}
       &\makebox[1em][l]{\( \upsilon\), \( \Upsilon \)} &upsilon（宇普西隆）  \pronounced{OOP-suh-LON}  \\
     \makebox[1em][l]{\( \iota \)} &iota（约塔） \pronounced{eye-OH-tuh}
       &\makebox[1em][l]{\( \phi \), \( \Phi \)} &phi（斐）  \pronounced{FEE, or FI (as in hi)}    \\
     \makebox[1em][l]{\( \kappa  \)} &kappa（卡帕）  \pronounced{KAP-uh}
       &\makebox[1em][l]{\( \chi \)}  &chi（凯）  \pronounced{KI (as in hi)}    \\
     \makebox[1em][l]{\( \lambda \), \( \Lambda \)} &lambda（拉姆达）  \pronounced{LAM-duh}
       &\makebox[1em][l]{\( \psi    \), \( \Psi \)}  &psi（普赛） \pronounced{SIGH, or PSIGH}    \\
     \makebox[1em][l]{\( \mu  \)}  &mu（缪）  \pronounced{MEW}
       &\makebox[1em][l]{\( \omega  \), \( \Omega \)} &omega（奥米伽）  \pronounced{oh-MAY-guh}
  \end{tabular}   \\[1.5ex]
  表中列出的大写字母是那些与相应拉丁字母大写写法不同的大写希腊字母。
\end{center}

\frontmatter
\chapter*{前言}
本书是吉姆·赫弗伦（Jim~Hef{}feron）《线性代数》
一书中全部习题的解答。
此处标为 One.II.3.4 的解答，
对应第一章、第二节、第三小节中编号为~4~的题目。
各章末尾的专题单独编号。

如果你拿到的是本文件的电子版，
请把它与主书放在同一目录下。
这样，在主书中点击题号即可跳到相应解答，
在这里点击解答编号即可跳回原题
（主书文件须名为 book-zh.pdf，本文件须名为 jhanswer-zh.pdf——
由于 PDF 链接语言的限制，这两个文件不能改名）。

欢迎勘误与建议。
联系方式见本书主页
\url{http://joshua.smcvt.edu/linearalgebra}。

\vspace{.5in}
\begin{raggedright}
Jim Hef{}feron     \\
Saint Michael's College, Colchester VT USA \\
\input{publicationdate}
\end{raggedright}
\clearemptydoublepage\tableofcontents
\setcounter{page}{1}
\mainmatter
\thispagestyle{empty}
\input{\ansfile}
\end{document}
